% =============================================================================
%  STYLE : couleurs, bannières vertes, en-tête et pied de page, titres, tables
% =============================================================================
\usepackage{etoc}

% --- Couleurs ----------------------------------------------------------------
\definecolor{vertnuit}{HTML}{1B4D1F}
\definecolor{vertfonce}{HTML}{2E7D32}
\definecolor{vertmoyen}{HTML}{66BB6A}
\definecolor{vertclair}{HTML}{DCEFD9}
\definecolor{bandeau}{HTML}{DDE1F0}
\definecolor{acompleter}{HTML}{E65100}

% --- Champ à compléter -------------------------------------------------------
\newcommand{\AC}[1]{\textcolor{acompleter}{[À COMPLÉTER : #1]}}

% --- Bannière verte (titres liminaires et intercalaires) ---------------------
\newcommand{\banniere}[1]{%
  \begin{center}
  \begin{tikzpicture}
    \node[draw=vertfonce, line width=1.2pt, rounded corners=10pt,
          top color=vertclair, bottom color=vertmoyen!70,
          text=vertnuit, font=\Large\bfseries, align=center,
          inner xsep=22pt, inner ysep=9pt, drop shadow={opacity=0.25}] (b) {#1};
    \foreach \s in {west,east}{
      \fill[vertfonce] (b.\s) circle (5pt);
      \fill[vertclair] (b.\s) circle (2.2pt);
    }
  \end{tikzpicture}
  \end{center}}

% --- Ruban vert du numéro de page --------------------------------------------
\newcommand{\ruban}[1]{%
  \tikz[baseline=(r.base)]{
    \node[fill=vertfonce, text=white, font=\small\bfseries,
          minimum width=1.2cm, inner ysep=3pt] (r) {#1};
    \fill[vertfonce] (r.north east) -- ++(0.35,0) -- ($(r.east)+(0.15,0)$)
                     -- ($(r.south east)+(0.35,0)$) -- (r.south east) -- cycle;
    \fill[vertfonce] (r.north west) -- ++(-0.35,0) -- ($(r.west)+(-0.15,0)$)
                     -- ($(r.south west)+(-0.35,0)$) -- (r.south west) -- cycle;
  }}

% --- En-tête et pied de page -------------------------------------------------
\newcommand{\enteteRapport}{%
  \colorbox{bandeau}{\parbox{\dimexpr\textwidth-2\fboxsep\relax}{\centering
    \footnotesize\bfseries\itshape \InfoTheme{} -- cas de \InfoEntrepriseSigle{}
    au chantier \InfoChantier}}}
\newcommand{\styleRapport}{%
  \fancyhf{}%
  \fancyhead[C]{\enteteRapport}%
  \fancyfoot[L]{\footnotesize\bfseries\itshape Rédigé et présenté par \InfoEtudiantCourt}%
  \fancyfoot[R]{\ruban{\thepage}}%
  \renewcommand{\headrulewidth}{0pt}%
  \renewcommand{\footrulewidth}{0.4pt}}
\pagestyle{fancy}\styleRapport
\fancypagestyle{plain}{\styleRapport}

% --- Numérotation : PARTIE > I. > 1. > a) ------------------------------------
\setcounter{secnumdepth}{3}
\setcounter{tocdepth}{3}
\renewcommand{\thesection}{\Roman{section}.}
\renewcommand{\thesubsection}{\arabic{subsection}.}
\renewcommand{\thesubsubsection}{\alph{subsubsection})}
\counterwithout{figure}{chapter}
\counterwithout{table}{chapter}
\counterwithout{equation}{chapter}
\renewcommand{\thefigure}{\arabic{figure}}
\renewcommand{\thetable}{\arabic{table}}
\renewcommand{\theequation}{\arabic{equation}}

\titleformat{\section}{\normalfont\large\bfseries\color{vertnuit}}{\thesection}{0.5em}{}
\titleformat{\subsection}{\normalfont\normalsize\bfseries\color{vertfonce}}{\thesubsection}{0.5em}{}
\titleformat{\subsubsection}{\normalfont\normalsize\bfseries\itshape}{\thesubsubsection}{0.5em}{}
\titlespacing*{\section}{0pt}{14pt}{6pt}
\titlespacing*{\subsection}{0pt}{10pt}{4pt}
\titlespacing*{\subsubsection}{0pt}{8pt}{3pt}

% --- Pages titrées non numérotées (liminaires, introduction, conclusion…) ----
\newcommand{\pageTitre}[1]{%
  \clearpage\phantomsection
  \addcontentsline{toc}{chapter}{#1}%
  \banniere{#1}\vspace{0.6em}}

% --- Intercalaire de partie --------------------------------------------------
\newcommand{\partie}[2]{%
  \clearpage\setcounter{section}{0}\phantomsection
  \addcontentsline{toc}{chapter}{#1 : #2}%
  \vspace*{\fill}
  \banniere{#1 :\\[4pt] #2}
  \vspace*{\fill}\clearpage}

\newcommand{\conclusionPartielle}{%
  \par\medskip\begin{center}\textbf{\underline{Conclusion partielle}}\end{center}}

% --- Tables des matières, figures, tableaux ----------------------------------
\renewcommand{\cftchapfont}{\bfseries}
\renewcommand{\cftchappagefont}{\bfseries}
\renewcommand{\cftchapleader}{\cftdotfill{\cftdotsep}}
\setlength{\cftbeforechapskip}{4pt}
\setlength{\cftsecnumwidth}{2.6em}
\setlength{\cftsubsecnumwidth}{2em}
\setlength{\cftsubsubsecnumwidth}{2em}
\renewcommand{\cftfigpresnum}{Figure~}
\renewcommand{\cftfigaftersnum}{~:}
\setlength{\cftfignumwidth}{5.6em}
\setlength{\cftfigindent}{0pt}
\renewcommand{\cfttabpresnum}{Tableau~}
\renewcommand{\cfttabaftersnum}{~:}
\setlength{\cfttabnumwidth}{6.2em}
\setlength{\cfttabindent}{0pt}
% Listes des figures et des tableaux : bannière + liste, sans \chapter* ni saut de page
\makeatletter
\AtBeginDocument{%
  \renewcommand{\listoffigures}{\banniere{LISTE DES FIGURES}\vspace{0.4em}\@starttoc{lof}}%
  \renewcommand{\listoftables}{\banniere{LISTE DES TABLEAUX}\vspace{0.4em}\@starttoc{lot}}}
\makeatother

% --- Tableaux ----------------------------------------------------------------
\newcolumntype{Y}{>{\raggedright\arraybackslash}X}
\newcommand{\entete}{\rowcolor{vertclair}}
\renewcommand{\arraystretch}{1.25}

% --- Listes ------------------------------------------------------------------
\setlist[itemize]{label=\textcolor{vertfonce}{\small$\blacklozenge$}, leftmargin=1.4cm, itemsep=2pt, topsep=3pt}
\setlist[enumerate]{leftmargin=1.4cm, itemsep=2pt, topsep=3pt}

% --- Figure factice ou réelle : \figureimg{fichier}{largeur}{légende}{label} --
\newcommand{\figureimg}[4]{%
  \begin{figure}[H]\centering
    \includegraphics[width=#2]{#1}
    \caption{#3}\label{#4}
  \end{figure}}
